\documentclass[11pt,letterpaper]{article}
\usepackage[spanish]{babel}
\usepackage[utf8]{inputenc}
\usepackage[T1]{fontenc}
\usepackage{geometry}
\usepackage{array}
\usepackage{booktabs}
\usepackage{hyperref}
\usepackage{xcolor}

\geometry{margin=1.8cm}
\hypersetup{colorlinks=true,linkcolor=blue,urlcolor=blue}
\setlength{\parskip}{0.45em}
\setlength{\parindent}{0pt}

\begin{document}

\section*{Resultados del Proyecto SAVD}

Repositorio: \href{https://github.com/Elnuget/Proyecto-savd}{https://github.com/Elnuget/Proyecto-savd}

El proyecto SAVD define y materializa una carcasa inicial para un Sistema de Automatizacion y Verificacion Documental. La propuesta parte de una planificacion agil con Scrum y se organiza en tres incrementos funcionales: gestion documental, automatizacion de validaciones y generacion de reportes con indicadores.

\subsection*{Resultado general}

Se construyo una base navegable del software esperado. Esta base permite demostrar el flujo principal del producto sin depender todavia de infraestructura externa, backend o base de datos. El prototipo permite iniciar sesion con credenciales de demostracion, cargar documentos permitidos, consultar registros, visualizar datos extraidos simulados, revisar estados de validacion, detectar observaciones y descargar un reporte basico.

\subsection*{Resultados por incremento}

\begin{tabular}{p{2.2cm}p{5.5cm}p{6.5cm}}
\toprule
\textbf{Sprint} & \textbf{Objetivo} & \textbf{Resultado obtenido} \\
\midrule
Sprint 1 & Gestion y almacenamiento de documentos & Interfaz base con autenticacion demostrativa, carga de archivos, validacion de formatos permitidos y listado de documentos registrados. \\
Sprint 2 & Automatizacion del proceso documental & Simulacion de extraccion de informacion, ejecucion de reglas basicas de validacion y asignacion de estado por documento. \\
Sprint 3 & Reportes y cierre del proceso & Panel con indicadores generales, detalle de resultado final, observaciones y descarga de reporte en texto. \\
\bottomrule
\end{tabular}

\subsection*{Historias de usuario cubiertas en la carcasa}

\begin{itemize}
  \item HU01: acceso mediante credenciales de demostracion.
  \item HU02: carga de documentos con formatos permitidos.
  \item HU03: listado actualizado de documentos cargados.
  \item HU04: visualizacion de informacion extraida simulada.
  \item HU05: validacion automatica mediante reglas basicas.
  \item HU06: seguimiento de estado del proceso por documento.
  \item HU07: consulta del resultado final de validacion.
  \item HU08: generacion y descarga de reporte basico.
  \item HU09: indicadores generales de documentos procesados.
\end{itemize}

\newpage

\section*{Resultados Tecnicos y Proyeccion}

\subsection*{Componentes entregados}

\begin{tabular}{p{4.2cm}p{10cm}}
\toprule
\textbf{Componente} & \textbf{Resultado} \\
\midrule
Interfaz principal & Pagina web estatica con secciones para acceso, carga, listado, resultados e indicadores. \\
Logica de aplicacion & Archivo JavaScript con control de sesion demostrativa, validacion de extensiones, generacion de registros, calculo de metricas y descarga de reportes. \\
Estilos visuales & Hoja CSS responsiva para uso en escritorio y dispositivos moviles. \\
Documentacion & README con descripcion del proyecto, alcance, ejecucion, estructura y proyeccion del software final. \\
Resultados academicos & Documento LaTeX de dos hojas sin portada con resumen de resultados del proyecto. \\
\bottomrule
\end{tabular}

\subsection*{Indicadores demostrados}

La carcasa muestra tres indicadores basicos: cantidad total de documentos cargados, documentos aprobados y documentos con observaciones. Estos indicadores representan el punto de partida para un tablero real de supervision documental.

\subsection*{Alcance pendiente para version productiva}

Para convertir la carcasa en un sistema completo se requiere implementar backend, base de datos, almacenamiento seguro de archivos, autenticacion real, motor OCR o extraccion de texto, reglas configurables de validacion, trazabilidad, auditoria, reportes PDF o Excel y pruebas automatizadas.

\subsection*{Conclusion}

El resultado obtenido cumple la funcion de transformar la planificacion Scrum en una base visible y ejecutable del sistema. La carcasa permite validar el flujo esperado con usuarios, alinear criterios de aceptacion y servir como punto de partida para el desarrollo incremental del producto final.

\end{document}
